\section{Обзор литературы}

В данной главе рассматриваются существующие алгоритмы и системы сжатия, поддерживающие произвольный доступ к данным, а также анализируются их ограничения. Цель обзора --- обосновать необходимость разработки новой библиотеки, сочетающей прозрачное сжатие с возможностью произвольной записи, вставки и удаления и возможностью подключения произвольного алгоритма сжатия.

\subsection{Классические алгоритмы сжатия}

Большинство популярных алгоритмов сжатия (например, \textsc{Deflate}, \textsc{Brotli}, \textsc{LZ4}, \textsc{Zstd}) работают как непрерывный поток кодов. В такой организации произвольный доступ к сжатым данным невозможен без полной распаковки всего потока с начала, что делает работу с большими файлами неэффективной.

Одним из распространённых решений является разбиение входных данных на блоки, каждый из которых сжимается независимо. Это позволяет распаковывать только нужный блок, не затрагивая остальные. Примером алгоритма, всегда разбивающего данные на независимые блоки, служит \texttt{bzip2}~\cite{bzip2}: каждый блок содержит собственную сжатую последовательность и сигнатуру начала, что даёт возможность произвольного чтения после построения индекса смещений. Аналогично, алгоритм \textsc{LZMA}~\cite{lzma} может работать в несплошном режиме (\textit{non-solid}), где блоки сжимаются независимо, обеспечивая ту же возможность.

Однако даже при блочной организации произвольная \textit{запись} (вставка, удаление или замена данных внутри блока) остаётся проблематичной. Изменение содержимого блока почти всегда меняет его сжатый размер, что приводит к сдвигу всех последующих блоков и требует перестройки их индексов и смещений. На практике любая модификация сжатых данных без полной перепаковки хотя бы затронутого блока и всех последующих оказывается невозможной. Поэтому для поддержки эффективной записи и редактирования требуются более сложные структуры данных и форматы хранения, выходящие за рамки классических алгоритмов.

\subsection{Универсальные системы с произвольным чтением}

Существуют надстройки, которые используют существующие алгоритмы сжатия и предоставляют произвольный доступ к данным без полной декомпрессии. Они реализуют дополнительные методы (индексацию, сохранение состояния декомпрессора и т.п.) поверх стандартных сжатых потоков.

\subsubsection{Zran}
Zran~\cite{zran} --- пример из дистрибутива \texttt{zlib}, добавляющий произвольное чтение к потоку Deflate. Он строит индекс из точек останова (\textit{break points}), для каждой из которых сохраняется полное состояние декомпрессора (скользящее окно, позиции в битовом потоке). При запросе на чтение распаковка начинается с ближайшей предыдущей точки останова, что требует декомпрессии лишь некоторого фрагмента. Индекс не требует модификации исходных сжатых данных и создаётся отдельно.

При формировании файла, поддерживающего произвольное чтение, сам индекс увеличивает общий размер архива. Накладные расходы зависят от частоты точек останова и объёма сохраняемого состояния: чем чаще точки, тем быстрее доступ, но тем больше индекс. Интервал между точками останова должен значительно превышать размер окна (по умолчанию в Deflate --- 32 КБ), иначе индекс становится чрезмерно большим. С другой стороны, слишком редкие точки останова приводят к необходимости распаковывать большие объёмы данных при каждом обращении, что замедляет доступ. 

Такой подход оправдан в сценариях, где запросы на чтение имеют объём, сопоставимый с интервалом между точками останова и существенно превышают размер окна. Таким образом, Zran ориентирован на крупные пакетные чтения, а интервал точек останова следует выбирать соответствующим образом.

\subsubsection{Seekable Zstd}

Seekable Zstd~\cite{zstd-seekable} --- надстройка над алгоритмом сжатия Zstd, добавляющая произвольное чтение. Исходные данные разбиваются на независимые кадры (frames), каждый из которых сжимается обычным Zstd с полным сбросом состояния. Для обеспечения произвольного доступа при формировании архива создаётся индекс, который хранит для каждого кадра его смещение в сжатом потоке и размер исходных данных. Индекс представляет собой простой статический массив, поиск по которому осуществляется с помощью бинарного поиска. При запросе на чтение по произвольной позиции библиотека по индексной таблице определяет кадры, содержащий нужные данные, и распаковывает только эти кадры.

Размер индексной таблицы пропорционален количеству кадров. Выбор размера кадра определяет компромисс: маленькие кадры дают более быстрый доступ к случайным позициям (распаковывается меньше данных), но снижают степень сжатия из-за потери межкадровой избыточности и увеличивают накладные расходы на индекс. Крупные кадры повышают степень сжатия, но требуют распаковки большего объёма данных при каждом обращении. Однако, в отличии от \textsc{Zran}, размер блока не обязан быть достаточно большим, так как размер индекса не растёт так сильно, что делает его пригодным для работы даже при маленьких операциях чтения.

Важно заметить, что хоть \textsc{seekable zstd} реализован конкретно для \textsc{Zstd}, подобный подход с разделением на блоки является вполне универсальным и может быть реализован с любым алгоритмом сжатия (примерами служат упомянутые выше \textsc{bzip2} и \textsc{LZMA}).

\subsection{Специализированные решения}

% TODO (CRAM, LCQS, FITS, ...?)